%\documentclass[minimal,compress, draft]{beamer} %draft pour pas mettre image et compil rapidement
\documentclass[minimal,compress]{beamer}
%\documentclass[minimal,compress]{thesis}


%virer la barre moche
\setbeamertemplate{navigation symbols}{}


%gets rid of bottom navigation bars
%\setbeamertemplate{footline}[page number]{}

\usepackage[utf8]{inputenc}
%\usepackage[english]{babel}
\usepackage[french]{babel}
\usepackage{hyperref}
\usepackage{graphicx}
\usepackage{eso-pic}	%put logo anywhere
\usepackage{xcolor}


\definecolor{mydarkred}{rgb}{0.40,0,0}
\definecolor{myred}{rgb}{0.50,0.17,0.17}
\definecolor{mydarkgray}{rgb}{0.33,0.33,0.33}
\definecolor{mygray}{rgb}{0.47,0.47,0.47}
\definecolor{mydarkblue}{rgb}{0.16,0.26,0.46}
\definecolor{myblue}{rgb}{0.43,0.50,0.64}
\definecolor{mydarkgreen}{rgb}{0.16,0.46,0.26}
\definecolor{mygreen}{rgb}{0.43,0.64,0.50}

\beamerboxesdeclarecolorscheme{redblock}{mydarkred}{myred!50!white}
\beamerboxesdeclarecolorscheme{greenblock}{mydarkgreen}{mygreen!50!white}

%http://mcclinews.free.fr/latex/beamergalerie/completsgalerie.html
%\usetheme{Dresden}
%\usetheme{Frankfurt}
%\usetheme{Montpellier}
\usetheme{monTheme}			%http://titilog.free.fr/index.htm


%entete pied page http://mcclinews.free.fr/latex/beamergalerie/outerdetails.html
%\useoutertheme{shadow}
%\useoutertheme[subsection=false]{Smoothbars}
%\useoutertheme{infolines}
%\useoutertheme{miniframes}
%\useoutertheme{smoothtree}
%\useoutertheme{tree}

%ensemble des sections dans la headline
%http://mcclinews.free.fr/latex/beamermodif/beamersection.html
\defbeamertemplate*{headline}{}
{%
  \begin{beamercolorbox}{section in head/foot}
    \insertsectionnavigationhorizontal{\paperwidth}{}{}
  \end{beamercolorbox}%
 }

%\defbeamertemplate*{headline}{miniframes theme}
%{%
%  \begin{beamercolorbox}{section in head/foot}
%    \insertsectionnavigationhorizontal{\paperwidth}{}{}
%  \end{beamercolorbox}%
%}

%shadow sous la healine depuis le beamerthemeshadow
%\setbeamercolor{frametitle}{parent=subsection in head/foot}
%\setbeamercolor{frametitle right}{parent=section in head/foot}
%\pgfdeclarehorizontalshading[frametitle.bg,frametitle right.bg]{beamer@frametitleshade}{\paperheight}{%
%  color(0pt)=(frametitle.bg);
%  color(\paperwidth)=(frametitle right.bg)}
%\AtBeginDocument{
%  \pgfdeclareverticalshading{beamer@topshade}{\paperwidth}{%
%    color(0pt)=(bg);
%    color(4pt)=(black!50!bg)}
%}
%\addtobeamertemplate{headline}
%{}
%{%
%  \vskip-0.2pt
%  \pgfuseshading{beamer@topshade}
%  \vskip-2pt
%}






\setbeamertemplate{frametitle}[default]
\setbeamertemplate{blocks}[rounded][shadow=true]

\usefonttheme{structurebold}
\setbeamerfont{section in toc}{series=\bfseries}
\setbeamerfont{alerted text}{series=\bfseries}
\setbeamerfont{quote}{series=\normalfont}

%\setbeamercolor{section in head}{fg=lightgray,bg=mydarkgray}
\setbeamercolor{section in head/foot} {fg=lightgray,bg=mydarkgray}
%\setbeamercolor{section in foot} %{fg=lightgray,bg=mydarkgray}
%{fg=white,bg=white}
\setbeamercolor{subsection in head/foot}{fg=mydarkred,bg=}
\setbeamercolor{titlelike}{fg=mydarkred,bg=}
\setbeamercolor{section in toc}{fg=mydarkred,bg=}
\setbeamercolor{subsection in toc}{fg=black!,bg=}
\setbeamercolor{item}{fg=mydarkred,bg=}
\setbeamercolor{subitem}{fg=mydarkgray,bg=}
\setbeamercolor{subsubitem}{fg=myred,bg=}
\setbeamercolor{navigation symbols}{fg=mygray,bg=}
\setbeamercolor{block title example}{fg=mydarkred,bg=}
\setbeamercolor{block title}{fg=mydarkred,bg=}
\setbeamercolor{block body}{fg=,bg=}
%\setbeamercolor{frametitle}{fg=mydarkgray,bg=}
\setbeamercolor{frametitle}{fg=mydarkred,bg=}
\setbeamercolor{alerted text}{fg=mydarkred,bg=}

%\pgfpagesuselayout{4 on 1}[a4paper,border shrink=5mm,landscape]

%In order to get a full page printable presentation (i.e. auto resized to A4 paper)
%\pgfpagesuselayout{resize to}[a4paper,border shrink=0mm,landscape]



%logo dauphine
\newcommand\AtPagemyUpperLeft[1]{\AtPageLowerLeft{%
\put(\LenToUnit{0.8\paperwidth},\LenToUnit{0.87\paperheight}){#1}}}
\AddToShipoutPictureFG{
  \AtPagemyUpperLeft{{\includegraphics[width=2cm,keepaspectratio]{logo}}}
}%


%numero pages
\addtobeamertemplate{footline}{\hspace{1em}\insertshorttitle\hfill\insertframenumber/\inserttotalframenumber\hspace{1em}\null}


%footer
%\makeatother
%\setbeamertemplate{footline}
%{
%  \leavevmode%
%  \hbox{%
%%  \begin{beamercolorbox}[wd=.4\paperwidth,ht=2.25ex,dp=1ex,center]{author in head/foot}%
%%    \usebeamerfont{author in head/foot}\insertshortauthor
%%  \end{beamercolorbox}%
%  \begin{beamercolorbox}[wd=.6\paperwidth,ht=2.25ex,dp=1ex,center]{title in head/foot}%
%    \usebeamerfont{title in head/foot}\insertshorttitle\hspace*{3em}
%    \insertframenumber{} / \inserttotalframenumber\hspace*{1ex}
%  \end{beamercolorbox}}%
%  \vskip0pt%
%}
%\makeatletter

%%%%%%%%%%%%%%%%%%%%%%%%%%%%%%%%%%%%%%%%%%%%%%%%%%%%%%%%%%%%%%%%%%%%%%%%%%%%%% 

\title[M1 MIAGE Apprentissage]
{MIAGE (Méthodes Informatiques Appliquées à la Gestion des Entreprises)}

\author{M1 - Formation en apprentissage}
\institute{}%
\date{}



%%%%%%%%%%%%%%%%%%%%%%%%%%%%%%%%%%%%%%%%%%%%%%%%%%%%%%%%%%%%%%%%%%%%%%%%%%%%%% 

\begin{document}



\frame[plain]{
 \titlepage
 
 
 \begin{tabular}{ll}
Responsable pédagogique Dauphine: & \textbf{Florian Sikora} \\
Chargée de mission CFA AFIA: & \textbf{Patricia Lavagna} \\
Administration: & \textbf{Delphine Galtier}
 \end{tabular}
}




\begin{frame}{Les disciplines}
\begin{itemize}
\item Technologies de l’\alert{informatique} et de la \alert{gestion}

\item Modélisation de \alert{systèmes d’information} et méthodes de \alert{conduite de projet}

\item Fonctionnement des organisations, leurs structures, et leurs impératifs stratégiques

\item L’anglais %{\scriptsize (plus une deuxième langue vivante)}

\item \alert{Professionnalisation} avec la formation par apprentissage
\end{itemize}
\end{frame}

\begin{frame}{Entreprises d'accueil}

\begin{itemize}
\item Banques, Assurances, Sociétés de Services, Startups, Logistique,...
\item Par exemple:
\end{itemize}

\centering
\begin{minipage}{.35\textwidth}
 Caisse des dépôts
 
 Dassault
 
 Air France
 
BNP Paribas

 EDF
 

\end{minipage}
\begin{minipage}{.35\textwidth}
 LCL
 
 Société Générale
 
GFI

 Natixis
 
 Canal+

\end{minipage}



\end{frame}

\begin{frame}{Organisation}
\begin{itemize}
\item Apprentissage
\begin{itemize}
\item 2 semaines entreprise / 2 semaines université
\item 	Maître d’apprentissage en entreprise + tuteur enseignant
\item 	Examen au fil de l’eau
\item 	Présence obligatoire cours/TD/TP

\end{itemize}
\item Promotion
\begin{itemize}
\item 25 à 30 apprentis
\item L3 apprentissage  en grande partie. Quelques places pour le L3 formation initiale et les externes à l’université Paris Dauphine
\end{itemize}
\end{itemize}
\end{frame}



\begin{frame}{Programme S1}
\begin{itemize}
\item \alert{Bloc Fondamental 1} Semestre 1 (24 ECTS)
\begin{itemize}
\item Systèmes de Gestion de Bases de Données  (4 ECTS)
\item Programmation  Java Avancée (4 ECTS)
\item Systèmes d’Information Avancés 1 (3 ECTS)
\item Intelligence Artificielle et raisonnement (4 ECTS)


\item Applications Réseaux (3 ECTS)
\item Analyse et fouille de données (3 ECTS)
\item Marketing (3 ECTS)	
\end{itemize}
\end{itemize}
\end{frame}

\begin{frame}{Programme S1}
\begin{itemize}
\item \alert{Bloc communication} Semestre 1 (6 ECTS)
\begin{itemize}
\item Organisation – Communication (3 ECTS)
\item Anglais 1 (3 ECTS)
%\item Deuxième langue vivante (3 ECTS)

\end{itemize}
\end{itemize}
\end{frame}




\begin{frame}{Programme S2}
\begin{itemize}
\item \alert{Bloc Fondamental 2} Semestre 2 (16 ECTS)
\begin{itemize}
\item Systèmes et Algorithmes Répartis  (4 ECTS)
\item Systèmes d’Information Avancés 2 (4 ECTS)
\item Analyse Financière (3 ECTS) 
\item e-applications (4 ECTS)
\item Anglais 2 (1 ECTS)

\end{itemize}
\end{itemize}
\end{frame}

\begin{frame}{Programme S2}
\begin{itemize}
\item \alert{Options} Semestre 2 (6 ECTS) % (une \textcolor{mydarkgreen}{verte} et une \textcolor{mydarkblue}{bleue})
\begin{itemize}
\item Programmation C++ (3 ECTS)
\item {Marchés financiers (3 ECTS)} 
\item {Optimisation combinatoire (3 ECTS)} 
\item {Ordonnancement et  gestion de production (3 ECTS)}
\item {Décision et Analyse Multicritère (3 ECTS)} 
\end{itemize}
\item 2 options à choisir.
\end{itemize}
\end{frame}








\begin{frame}{Programme}

\begin{itemize}
\item \alert{Bloc entreprise} (8 ECTS)
\end{itemize}

\begin{itemize}
\item Modalités:
\begin{itemize}
\item Mémoire sur une problématique rencontrée dans le cadre de la mission
\item Soutenance (première semaine de septembre).
\end{itemize}
\item Évaluation:
\begin{itemize}
\item Pertinence du sujet
\item Qualité du travail réalisé (recherche, analyse et synthèse)
\item Qualité de la rédaction du mémoire
\item Travail en entreprise
\end{itemize}

\end{itemize}

\end{frame}


\begin{frame}{Contrôle des connaissances}
\begin{itemize}
\item Validation de l’année:
\begin{itemize}
\item Note finale de \alert{chacun des blocs} $\geq$ \alert{10}/20
\item Note finale de \alert{chaque UE} $\geq$ \alert{5}/20.
\item Valider 60 ECTS.
\end{itemize}
\item Seconde session:
\begin{itemize}
\item Repasser les UEs avec une note $<$ \alert{5/20}
\item Compensation des UE dans un bloc
\item Choix de repasser ou non une UE avec note finale dans [5,10[
\end{itemize}
\end{itemize}
\end{frame}

\begin{frame}{M2 MIAGE en apprentissage}

\begin{itemize}
\item Accès à la \alert{2ème année} du Master Informatique des Organisations.
\item Spécialités:
\begin{itemize}
\item \alert{MIAGE-IF} (Informatique pour la finance):
\begin{itemize}
\item Forme des informaticiens dans les métiers de la \alert{banque et de l’assurance}
\end{itemize}
\item \alert{MIAGE-SITN} (Sys. d’info. et tecno. nouv.)
\begin{itemize}
\item Forme des professionnels en conception et mise en œuvre de \alert{systèmes d’information}
\end{itemize}
\item \alert{ID-MIAGE} (Informatique Décisionnelle)
\begin{itemize}
\item Forme des informaticiens en \alert{aide à la décision et en Informatique Décisionnelle}
\end{itemize}
\end{itemize}
\end{itemize}

\end{frame}





%%%%%%%%%%%%%%%%%%%%%%%%%%%%%%%%%%%%%%%%%%%%%%%%%%%%%%%%%%%%%%%%%%%%%%%%%%%%%% 

\end{document}